% ============================================================================
%  packages.tex — pacchetti. Nessuna definizione di stile qui.
%  L'ordine conta: hyperref va caricato quasi per ultimo, cleveref dopo hyperref.
% ============================================================================

% --- Lingua e codifica -------------------------------------------------------
\usepackage[T1]{fontenc}
\usepackage[utf8]{inputenc}
\usepackage[italian]{babel}
\usepackage{csquotes}          % obbligatorio con biblatex; fornisce \enquote{}

% --- Layout ------------------------------------------------------------------
\usepackage[a4paper,margin=2.5cm,bindingoffset=0.5cm]{geometry}
\usepackage{emptypage}         % pagine bianche davvero vuote (niente header)
\usepackage{fancyhdr}
\usepackage{setspace}

% --- Matematica --------------------------------------------------------------
\usepackage{amsmath}
\usepackage{amssymb}
\usepackage{amsthm}
\usepackage{mathtools}
\usepackage{bm}

% --- Figure e tabelle --------------------------------------------------------
\usepackage{graphicx}
\usepackage{float}
\usepackage{subcaption}
\usepackage{booktabs}          % toprule/midrule/bottomrule: mai \hline
\usepackage{tabularx}
\usepackage{longtable}
\usepackage{multirow}
\usepackage{caption}
\captionsetup{font=small,labelfont=bf,width=0.9\linewidth}

% --- Codice ------------------------------------------------------------------
\usepackage{listings}
\usepackage{xcolor}

% --- Bibliografia ------------------------------------------------------------
%  authoryear: "(Elgammal et al., 2017)". Per passare al numerico IEEE basta
%  sostituire questa riga con: \usepackage[style=ieee,backend=biber]{biblatex}
\usepackage[
  style=authoryear,
  backend=biber,
  maxcitenames=2,
  maxbibnames=99,
  giveninits=true,
  uniquename=init,
  sorting=nyt,
  doi=true,
  url=true,
  isbn=false
]{biblatex}
\addbibresource{references/bibliography.bib}

% --- Varie -------------------------------------------------------------------
\usepackage{enumitem}
\usepackage{xurl}              % spezza gli URL lunghi
\usepackage{xspace}
\usepackage{todonotes}
\usepackage[acronym,nomain,nonumberlist]{glossaries}

% --- Collegamenti (penultimo) ------------------------------------------------
\usepackage{hyperref}
\hypersetup{
  colorlinks=true,
  linkcolor=black,
  citecolor=black,
  urlcolor=blue!60!black,
  pdftitle={\tesiTitolo},
  pdfauthor={\tesiAutore}
}

% --- Riferimenti intelligenti (ultimo, dopo hyperref) ------------------------
\usepackage[italian,capitalise,noabbrev]{cleveref}
